\typeout{NT FILE chapter9.tex}

\chapter{Evaluation and Validation}
\label{cha:evaluation}

This chapter summarizes how the finite-state transducer (FST) extensions were validated in both the OCamlFLAT library and
the OFLAT application. It covers test layers, manual verification flows, and representative use cases that exercise some of the
most relevant features.

\section{Testing Methodology}
\label{sec:eval_methodology}

Validation was performed in three layers:
\begin{itemize}
  \item \textbf{Library unit tests (OCamlFLAT):} Tests live alongside model files
    (e.g., \texttt{TransducerTests.ml}) and target most functions such as \texttt{acceptOut}, \texttt{isDeterministic},
    and \texttt{minimize}.
  \item \textbf{Integration checks (OFLAT):} The JS build produced by \texttt{js\_of\_ocaml} was tested in the
  browser: loading the app, creating/editing FST edges with input:output labels, running step-by-step execution, and
  verifying that the info box and poppers update with accumulated output and configuration counts. Exercises using the
      \texttt{input->output} format were validated to confirm UI feedback matches library behavior. Minimization and
  determinization actions were also triggered from the UI to ensure results matched library functions.
  \item \textbf{Manual exploratory runs:} For edge cases (epsilon outputs, nondeterministic branches, conflicting outputs,
    and invalid output symbols), transducers were crafted in the UI and run through the stepper to confirm highlighting,
    table rebuilds, and error surfacing behaved as described in Chapters~\ref{cha:oflat-design} and \ref{cha:oflat-implementation}.
\end{itemize}

\paragraph{Stress and scale checks.}
We exercised transducers with high branching factors to observe exponential blow-up in trails and UI responsiveness, and ran determinization on conflicting-output FSTs to ensure non-termination is detected and reported as “undefined/failed” rather than hanging the UI.

\section{Use Cases}
\label{sec:eval_usecases}

Each use case was exercised end-to-end: constructing the model in OFLAT, stepping through execution, and comparing results
with the OCamlFLAT library output. This is not an exhaustive list of all features.

\subsection{Case 1: Simulating a Moore Machine}
\label{sec:eval_case_moore}
A deterministic Moore machine was created with uniform outputs per state. The UI was used to:
\begin{itemize}
  \item Draw the graph and label edges with input symbols (output fixed by state), then run the stepper to confirm the info
    box shows output determined solely by the current state.
  \item Run the exercise checker with \texttt{inside} examples (e.g., \texttt{a->x}, \texttt{ab->xy}) to verify
    	exttt{acceptExpectedOutput} passes and that the trace table displays the expected output word.
  \item Invoke the property panel to confirm the \texttt{moore} and \texttt{deterministic} flags are satisfied.
\end{itemize}

\begin{figure}[H]
  \centering
  \includegraphics[width=0.95\textwidth]{mooretrace.png} 
	\caption{Moore machine execution highlighting state-based outputs and trace table.}
  \label{fig:eval-moore}
\end{figure}

\subsection{Case 2: Non-Deterministic FST Execution}
\label{sec:eval_case_ndfst}
A nondeterministic FST with multiple outgoing edges per input was used to validate branching visuals and output tracking:
\begin{itemize}
	\item Step-by-step execution highlighted all active states; the chosen accepting path was overpainted, while alternative
		branches stayed visible. The info box showed different accumulated outputs per branch.
	\item Poppers reflected configuration multiplicity; advancing a step recomputed counts without rebuilding the graph.
\end{itemize}

\begin{figure}[H]
	\centering
  \includegraphics[width=0.95\textwidth]{fst_poppers.png} 
	\caption{Nondeterministic FST execution with branch highlights, popper counts, and output tracking.}
	\label{fig:eval-ndfst}
\end{figure}

\subsection{Case 3: FST Minimization}
\label{sec:eval_case_min}
Minimization was tested on a deterministic FST with redundant states:
\begin{itemize}
  \item The minimization routine produced a reduced graph; the UI recolored equivalence classes and redrew edges using
    \texttt{inputEdges} to preserve input:output labels.
  \item Execution before and after minimization yielded identical input/output behaviour (checked via \texttt{acceptOut}),
    and the trace table matched across versions.
  \item Property checks displays that both \texttt{deterministic} and \texttt{minimal} properties were satisfied after the minimization.
\end{itemize}

\begin{figure}[H]
  \centering
  \includegraphics[width=0.95\textwidth]{fst_rendermin.png} 
  \caption{FST minimization view with equivalence-class coloring and preserved input:output labels.}
  \label{fig:eval-min}
\end{figure}

\subsection{Case 4: Determinization Failure (Conflicting Outputs)}
\label{sec:eval_case_detfail}
An FST with conflicting outputs for the same (state, input),
and one with epsilon transitions leading to nondeterminism was tested:
Attempting determinization surfaced library errors, 
and the UI displayed the same message and preserved the original graph.
\begin{figure}[H]
  \centering
  \includegraphics[width=0.95\textwidth]{det_failure1.png} 
  \caption{Determinization failure due to epsilon transitions; error surfaced while graph is preserved.}
  \label{fig:eval-detfail}
\end{figure}
\begin{figure}[H]
  \centering
  \includegraphics[width=0.95\textwidth]{det_failure2.png} 
  \caption{Determinization failure due to conflicting outputs; error surfaced while graph is preserved.}
  \label{fig:eval-detfail2}
\end{figure}

\subsection{Case 5: Output-Aware Exercises}
\label{sec:eval_case_exercises}
Exercises using the arrow notation were validated end to end:
\begin{itemize}
  \item Inside examples like \texttt{ab->xy} passed only when the produced output matched exactly; outside examples like \texttt{ab->xz} 
  were rejected and shown with the actual output in the feedback panel.
  \item The same cases run through the library's \texttt{acceptExpectedOutput} yielded identical verdicts, confirming consistency between UI and core logic.
  \item Invalid output symbols (not in the declared output alphabet) were flagged immediately, both in the UI and via \texttt{acceptCheckOutput}.
\end{itemize}
\begin{figure}[H]
  \centering
  \includegraphics[width=0.95\textwidth]{ex_feedback.png} 
  \caption{Output-aware exercise feedback with arrow notation, highlighting mismatches.}
  \label{fig:eval-exercises}
\end{figure}

\subsection{Case 6: Scale (Branching Explosion)}
\label{sec:eval_case_scale}
A highly branching nondeterministic FST was run to gauge responsiveness:
\begin{itemize}
  \item Trail sizes grew exponentially with input length; the stepper remained responsive for short inputs but slowed as branches multiplied.
  \item Poppers and the info box updated correctly with large configuration sets.
\end{itemize}
\begin{figure}[H]
  \centering
  \includegraphics[width=0.95\textwidth]{limitfst.png} 
  \caption{Scale test with high branching; UI remains responsive while showing many active configurations.}
  \label{fig:eval-scale}
\end{figure}

\section{Discussion and Limitations}
\label{sec:eval_discussion}
The graphical UI was not formally usability-tested; uniformity with existing models was prioritized, and a full UX study remains future work. 
Determinization non-termination is detected and reported as an undefined/failed result to avoid UI hangs. 
Performance degrades with exponential branching as expected; exploratory stress tests help bound practical input sizes.